\section{Entorno y datos de las mediciones}
\label{sec:apendice_mediciones}

Este apartado reúne el entorno, las versiones, las órdenes y los registros completos de las
mediciones descritas en la Sección~\ref{sec:metodologia}, con el fin de permitir su reproducción y
su contraste con los valores derivados de la Sección~\ref{sec:resultados}.

\subsection{Entorno y versiones}
\label{sec:apendice_entorno}

La Tabla~\ref{tab:apendice_entorno} resume los equipos y las herramientas empleados en cada grupo
de mediciones. Las etapas de generación de ambas herramientas se midieron en el mismo servidor,
mientras que las mediciones de gas se realizaron en redes locales, con comprobaciones adicionales
sobre Optimism Sepolia en el caso de SP1.

\begin{table}[htbp]
	\centering
	\caption{Equipos, redes y versiones de las herramientas empleadas en las mediciones.}
	\label{tab:apendice_entorno}
	\footnotesize
	\begin{tabular}{>{\raggedright\arraybackslash}p{3.3cm} >{\raggedright\arraybackslash}p{5.9cm} >{\raggedright\arraybackslash}p{5.9cm}}
		\toprule
		\textbf{Elemento} & \textbf{Noir} & \textbf{SP1} \\
		\midrule
		\multicolumn{3}{l}{\textit{Generación de las pruebas}} \\
		\addlinespace
		Equipo & \multicolumn{2}{p{11.8cm}}{Máquina virtual de Oracle Cloud Infrastructure, forma
		\texttt{VM.Standard3.Flex}, con 7 OCPU (14 hilos), 90 GB de memoria y 150 GB de
		almacenamiento para el área de intercambio, sistema operativo Linux y generación
		exclusivamente en el procesador.} \\
		\addlinespace
		Herramientas & Nargo 1.0.0-beta.22 (commit \texttt{c57152f9}) y Barretenberg
		5.0.0-nightly.20260522 & \texttt{sp1-sdk} 6.1.0, cadena de compilación de Succinct basada en
		\texttt{rustc} 1.94.0 y \texttt{SP1\_PROVER=cpu} \\
		\addlinespace
		Programa medido & \texttt{noir\_bls12\_381\_validator}, versión anterior del circuito de consenso & \texttt{sync-committee-program},
		66\,378\,401 ciclos \\
		\addlinespace
		Fecha & Última ejecución el 10 de agosto de 2026 & 21 y 22 de julio de 2026 \\
		\midrule
		\multicolumn{3}{l}{\textit{Verificación en la cadena de destino}} \\
		\addlinespace
		Equipo y red & Computador Mac con 18 GB de memoria, Anvil con la bifurcación BPO1 (reglas
		posteriores a Osaka) y red de Hardhat en el propio proceso (reglas de Prague) & Anvil 1.3.5 (commit \texttt{9979a41b}), bifurcación Cancun, identificador de
		cadena 31337, y Node 20.20.2 \\
		\addlinespace
		Compilador & solc 0.8.28, optimizador activado, \texttt{runs = 1} &
		solc 0.8.24, optimizador activado, \texttt{runs = 1}, \texttt{viaIR} desactivado, EVM
		\texttt{shanghai} \\
		\addlinespace
		Verificador & Un contrato por circuito, generado por \texttt{bb write\_solidity\_verifier -t evm}
		con Noir 1.0.0-beta.22 y Barretenberg 5.0.0-nightly.20260522 &
		\texttt{SP1VerifierGroth16} v6.1.0 de \texttt{sp1-contracts}, commit
		\texttt{d3629729}~\cite{sp1contracts} \\
		\addlinespace
		Red pública & Ethereum Sepolia, solo como cadena de origen de la ejecución de punta a punta & Optimism Sepolia, identificador 11155420, bloque
		51\,307\,662, verificador desplegado en
		\texttt{0xb69f2584CBcFf99a}\allowbreak\texttt{58C4e7002E8b89Af54a6f4e2} \\
		\addlinespace
		Fecha & Noviembre de 2026 & 11 de noviembre de 2026 \\
		\bottomrule
	\end{tabular}
\end{table}

\subsection{Órdenes ejecutadas}
\label{sec:apendice_ordenes}

Cada etapa de generación se ejecutó mediante \texttt{/usr/bin/time -v bash -c} seguido de la
orden correspondiente, que se transcribe a continuación tal como figura en los registros.

\begin{lstlisting}[basicstyle=\ttfamily\footnotesize, breaklines=true, columns=fullflexible]
# Noir
nargo compile
nargo execute witness
bb write_vk -b target/noir_bls12_381_validator.json -o target/
bb prove -b target/noir_bls12_381_validator.json -w target/witness.gz -k target/vk -o target/
bb verify -k target/vk -p target/proof -i target/public_inputs

# SP1 (SP1_PROVER=cpu)
cd program && cargo prove build
cargo run --release -p sync-committee-script --bin sync-committee -- --execute --data data.json
cargo run --release -p sync-committee-script --bin vkey
cargo run --release -p sync-committee-script --bin sync-committee -- --prove --data data.json
\end{lstlisting}

\subsection{Registros de las etapas de generación}
\label{sec:apendice_registros}

Las Tablas~\ref{tab:apendice_registros_noir} y~\ref{tab:apendice_registros_sp1} transcriben los
campos de \texttt{/usr/bin/time -v} de cada etapa en sus unidades originales. La memoria residente
máxima se expresa en kibibytes, y las entradas y salidas del sistema de archivos, en bloques de 512
bytes. Todas las etapas terminaron con el código de salida cero.

\begin{table}[htbp]
	\centering
	\caption{Registros de \texttt{/usr/bin/time -v} de las etapas de Noir.}
	\label{tab:apendice_registros_noir}
	\footnotesize
	\setlength{\tabcolsep}{4pt}
	\begin{tabular}{>{\raggedright\arraybackslash}p{3.1cm} *{5}{>{\raggedleft\arraybackslash}p{2.2cm}}}
		\toprule
		\textbf{Campo} & \textbf{Compilación} & \textbf{Testigo} & \textbf{Clave} & \textbf{Prueba} &
		\textbf{Verificación} \\
		\midrule
		Tiempo de reloj & 1:28:17 & 1:35,31 & 0:59,54 & 1:24,36 & 0:00,01 \\
		Tiempo de usuario (s) & 3220,38 & 92,06 & 236,39 & 623,91 & 0,01 \\
		Tiempo de sistema (s) & 161,42 & 3,35 & 9,33 & 19,99 & 0,00 \\
		Uso del procesador & 63\,\% & 100\,\% & 412\,\% & 763\,\% & 160\,\% \\
		Memoria residente máxima (KiB) & 90\,948\,544 & 4\,600\,272 & 18\,063\,240 & 19\,972\,232 &
		8\,768 \\
		Fallos de página mayores & 3\,043\,192 & 0 & 18 & 2 & 1 \\
		Entradas del sistema de archivos & 127\,531\,904 & 8 & 2\,148\,672 & 1\,472 & 48 \\
		Salidas del sistema de archivos & 223\,624 & 65\,344 & 16 & 40 & 0 \\
		\bottomrule
	\end{tabular}
\end{table}

\begin{table}[htbp]
	\centering
	\caption{Registros de \texttt{/usr/bin/time -v} de las etapas de SP1.}
	\label{tab:apendice_registros_sp1}
	\footnotesize
	\setlength{\tabcolsep}{4pt}
	\begin{tabular}{>{\raggedright\arraybackslash}p{3.6cm} *{4}{>{\raggedleft\arraybackslash}p{2.4cm}}}
		\toprule
		\textbf{Campo} & \textbf{Compilación} & \textbf{Ejecución} & \textbf{Clave} &
		\textbf{Prueba} \\
		\midrule
		Tiempo de reloj & 0:00,58 & 0:36,67 & 0:05,55 & 17:09,45 \\
		Tiempo de usuario (s) & 0,41 & 353,90 & 38,61 & 11\,429,80 \\
		Tiempo de sistema (s) & 0,17 & 3,91 & 0,45 & 386,92 \\
		Uso del procesador & 101\,\% & 975\,\% & 703\,\% & 1147\,\% \\
		Memoria residente máxima (KiB) & 122\,372 & 7\,080\,364 & 604\,972 & 90\,644\,680 \\
		Fallos de página mayores & 0 & 0 & 79 & 527\,632 \\
		Entradas del sistema de archivos & 0 & 0 & 201\,296 & 5\,899\,464 \\
		Salidas del sistema de archivos & 0 & 8 & 416 & 8 \\
		\bottomrule
	\end{tabular}
\end{table}

\subsection{Registros de las mediciones de gas}
\label{sec:apendice_gas}

La Tabla~\ref{tab:apendice_gas_sp1} detalla las mediciones de gas de SP1 sobre Anvil, con la
estimación obtenida mediante \texttt{eth\_estimateGas} y el consumo registrado en el recibo de
cada transacción. Ambas columnas incluyen el costo intrínseco de 21\,000 unidades y el de los datos de la transacción. En la red local, la
dirección del verificador fue \texttt{0x5FbDB2315678afecb367f032d93F642f64180aa3}, y su función
\texttt{VERIFIER\_HASH} devolvió
\texttt{0x4388a21c687fdd5f218d}\allowbreak\texttt{7e3d13190cac4c5355818d}\allowbreak\texttt{3605fd5fb811df468ee696}, el mismo valor que el
verificador desplegado en Optimism Sepolia. La prueba alterada, obtenida modificando el último byte
de la prueba, revirtió con el error \texttt{ProofInvalid()}, de selector \texttt{0x7fcdd1f4}, en
ambas redes. En Optimism Sepolia, el verificador desplegado devolvió una estimación de 259\,224
unidades de gas para la prueba válida, y la repetición de estos controles en el bloque
51\,361\,939 arrojó la misma estimación y confirmó los precios de los precompilados. Las
mediciones de \texttt{updateHeader} emplearon los valores públicos de una actualización real de
Ethereum Sepolia con 509 de 512 firmantes. Su código no es idéntico al compilado localmente, por lo que
esta cifra no se combina con las de la tabla.

\begin{table}[htbp]
	\centering
	\caption{Gas de las operaciones de SP1 medidas sobre Anvil.}
	\label{tab:apendice_gas_sp1}
	\footnotesize
	\begin{tabular}{>{\raggedright\arraybackslash}p{6.8cm} >{\raggedleft\arraybackslash}p{2.6cm} >{\raggedleft\arraybackslash}p{2.2cm} >{\raggedleft\arraybackslash}p{1.8cm}}
		\toprule
		\textbf{Operación} & \textbf{Estimación} & \textbf{Recibo} &
		\textbf{Datos (bytes)} \\
		\midrule
		\texttt{verifyProof} con la prueba Groth16 válida & 259\,608 & 253\,395 & 1\,156 \\
		Hash de los valores públicos de la prueba empleada & 29\,072 & 29\,072 & 676 \\
		Hash de valores públicos de 320 bytes, versión anterior del programa & 25\,825 & 25\,825 & 388 \\
		\texttt{updateHeader} con verificador simulado, primer anclaje & 85\,340 & 85\,340 & No informado \\
		\texttt{updateHeader} con verificador simulado, anclaje ordinario & 68\,240 & 68\,240 & No informado \\
		\texttt{setCommittee}, rotación del compromiso & 30\,127 & 30\,127 & 36 \\
		\texttt{updateHeader} con verificador simulado, tras la rotación & 68\,240 & 68\,240 & No informado \\
		\bottomrule
	\end{tabular}
\end{table}

La Tabla~\ref{tab:apendice_precompilados} recoge el gas mínimo con que cada contrato precompilado
de la curva BN254 produjo el resultado esperado en Optimism Sepolia. En todos los casos, la misma
llamada con una unidad de gas menos falló.

\begin{table}[htbp]
	\centering
	\caption{Gas de los contratos precompilados de BN254 comprobado en Optimism Sepolia.}
	\label{tab:apendice_precompilados}
	\footnotesize
	\begin{tabular}{>{\raggedright\arraybackslash}p{2.2cm} >{\raggedright\arraybackslash}p{6.6cm} >{\raggedleft\arraybackslash}p{2.6cm}}
		\toprule
		\textbf{Dirección} & \textbf{Operación} & \textbf{Gas} \\
		\midrule
		\texttt{0x06} & Suma de un punto consigo mismo & 150 \\
		\texttt{0x07} & Multiplicación del generador por dos & 6\,000 \\
		\texttt{0x08} & Emparejamiento de dos pares que se cancelan & 113\,000 \\
		\texttt{0x08} & Emparejamiento de cuatro pares que se cancelan & 181\,000 \\
		\bottomrule
	\end{tabular}
\end{table}

La Tabla~\ref{tab:apendice_circuitos_noir} detalla las características de los circuitos de Noir
que integran el puente y de sus contratos verificadores. En la ejecución de punta a punta, la
transacción de bloqueo en Ethereum Sepolia fue
\texttt{0xcc943f29e0b5973be022}\allowbreak\texttt{92d2f9f4e0615a49bd584d}\allowbreak\texttt{9e7a784f363cd4df8a73d3},
incluida en la posición 97 del bloque 12\,122\,052, y la firma del comité correspondió al
\textit{slot} 11\,565\,990, con 509 de 512 participantes. Las pruebas automatizadas de los
contratos emplearon, además, pruebas reales de otra firma de Ethereum Sepolia, del \textit{slot}
11\,565\,917, con 500 de 512 participantes, correspondiente al bloque 12\,121\,982, y las del
circuito de inclusión, transacciones y recibos reales de un bloque de una red local de Hardhat con
20 transacciones adicionales. El circuito de inclusión limita la transacción serializada a
256 bytes, frente a los 173 a 181 que ocupa una transacción de bloqueo, y con las cotas del árbol
de recibos alcanzaba 4,94 millones de compuertas, lo que lo habría llevado a $2^{23}$ filas.

\begin{table}[htbp]
	\centering
	\caption{Características de los circuitos de Noir que integran el puente.}
	\label{tab:apendice_circuitos_noir}
	\footnotesize
	\begin{tabular}{>{\raggedright\arraybackslash}p{4.2cm} *{3}{>{\raggedleft\arraybackslash}p{3.0cm}}}
		\toprule
		\textbf{Característica} & \textbf{Comité} & \textbf{Emparejamiento} & \textbf{Inclusión} \\
		\midrule
		Compuertas & 3\,117\,372 & 8\,170\,394 & 4\,105\,601 \\
		Tamaño del circuito (filas) & $2^{22}$ & $2^{23}$ & $2^{22}$ \\
		Tamaño de la prueba (bytes) & 11\,072 & 11\,456 & 11\,072 \\
		Entradas públicas & 7 & 1 & 168 \\
		Código del verificador (bytes) & 17\,532 & 17\,155 & 17\,407 \\
		Gas de la verificación, Osaka & 4\,489\,878 & 4\,578\,014 & 4\,745\,084 \\
		Gas de la verificación, Prague & 3\,077\,258 & 3\,133\,115 & 3\,332\,963 \\
		Tiempo de generación de la prueba & 23 s & 75 s & 23 s \\
		Tiempo de compilación & 78 s & 32 min & No informado \\
		\bottomrule
	\end{tabular}

	\vspace{0.5em}
	\begin{minipage}{0.95\textwidth}
		\footnotesize
		Los tiempos son aproximados y se obtuvieron en el computador Mac de 18 GB de memoria. Las
		siete entradas del circuito del comité son el compromiso del comité, las dos mitades de
		128 bits de la raíz de recibos y de la raíz de transacciones, el \textit{slot} y el valor de enlace, y la única del
		circuito de emparejamiento es ese mismo valor de enlace.
	\end{minipage}
\end{table}
